\section{真实对话、局限与未来方向}

最后一组 slide 把抽象模块重新落回对话。案例展示 Cicero 怎样用地图上可执行的替代方案说服玩家；限制 slide 则明确列出 intent 表示、truthfulness 和 domain-specific planning 尚未解决的问题。理解这些限制，才能避免把模块化成功误读成通用社会智能。

\subsection{奥地利案例：语言必须指向合法协作}

前面已经分别解释 intent、planner 和 filter，现在要检查它们是否在一段真实对话中共同落地。本节不按文采评价消息，而按“规则合法、利益可接受、失败可回退”三层标准阅读。slide 39 展示 Cicero 作为 Austria 与其他玩家的交互；地图圈出争夺区域，右侧对话把单位、位置和建议动作绑定起来。读图时还要检查对方若不合作，Cicero 是否仍有备选计划。

\lecturefigure{slide-39-dialogue-austria.jpg}{Cicero as Austria：以局面和行动为基础的协商}{官方录像恢复幻灯片 V039；Stanford CS25 Lecture 14}

这类消息之所以优于泛泛“我们结盟吧”，是因为它降低协调歧义。对方可以根据具体 orders 判断收益，Cicero 也能在下一轮比较承诺与实际行动，从而更新信任。

\subsection{法国与 Tunis：策略性说服而非漂亮续写}

slide 40 的 Tunis 对话更清楚地显示 persuasion。Cicero 不是只要求 Turkey 让步，而是提出能让 Turkey 获得其他收益的替代动作，同时为 France 腾出 Tunis。消息的质量来自对联合行动空间的重新构造：给对方一个可接受方案，才能改变其 policy。

\lecturefigure{slide-40-dialogue-france.jpg}{Cicero as France：围绕 Tunis 提出兼顾对方利益的替代方案}{官方录像恢复幻灯片 V040；Stanford CS25 Lecture 14}

\teachervoice{老师把这个例子称为 grounded strategic persuasion：语言不是独立“很有说服力”，而是背后确实存在一个对双方短期收益都说得通的行动组合。没有 planner，decoder 很容易生成听起来合作、实际却不可执行的建议。}

\subsection{四类局限}

limitations slide 把系统边界集中列出。第一，intent 主要表示近期 actions，难以表达长期目标、语气或关系策略；第二，dialogue model 与 strategy model 分开，前者不承担完整 world model；第三，value model 不读取 dialogue，因此无法估计撒谎造成的长期信任损失；第四，planning 仍是 Diplomacy-specific，并非讲者期待的通用推理机制。

\lecturefigure{slide-41-limitations.jpg}{Cicero 的限制与未来方向：intent、世界模型、信任和通用规划}{官方录像恢复幻灯片 V041；Stanford CS25 Lecture 14}

\begin{longtable}{L{0.22\textwidth}L{0.31\textwidth}L{0.37\textwidth}}
\toprule
限制 & 现有设计 & 直接后果 \\
\midrule
Intent 过窄 & 以近程 orders 为主要接口 & 难表达长期 alliance、交流风格和多步承诺 \\
世界模型分离 & planner 负责状态与价值，LM 负责表达 & 提高可控性，但跨模块错误可能累积 \\
Value model dialogue-blind & 不直接建模话语与长期关系 & 无法给“失去信任”可靠定价 \\
规划器领域化 & 利用 Diplomacy 规则与动作表示 & 不能直接迁移到开放世界任务 \\
文本模态有限 & 只处理 text-based communication & 难覆盖语音、视觉、时间压力与非语言信号 \\
\bottomrule
\end{longtable}

\subsection{Truthfulness 是工程约束，不是道德证明}

因为 value model 看不到 dialogue 的长期后果，它可能认为一条短期有效的谎言价值很高，却无法计算未来所有玩家不再信任 Cicero 的损失。团队因此让 dialogue model 从 truthful intents 出发，并过滤疑似不真实训练 turn。这是一种在模型盲区前收缩 action space 的安全设计。

\begin{warningbox}{不要从 Cicero 推出“真实总是最优”}
课堂证据只能支持：在当前表示和 evaluator 下，系统没有能力可靠评估长期欺骗，因此 truthfulness 是更稳健的工程选择。它没有证明所有 Diplomacy 局面都不需要 deception，也没有证明模型拥有稳定道德意图。
\end{warningbox}

\teachervoice{讲者非常坦率地说，truthful intention 的原因之一正是 value model 不理解长期 trust damage。这个解释比“我们让模型成为诚实的人”更精确，也展示了系统安全常见路径：当 evaluator 不会评估某类风险时，先限制生成空间。}

\subsection{从游戏走向 general planning}

closing slide 回顾 Cicero 的结果并给出论文、代码与模型入口。更重要的是 Q\&A：讲者并不期待再找一个“更难的桌游”重复同一路线，而希望研究能够作用于 foundation model 的通用 planning；文本输入输出是当时可行选择，却不是最终多模态智能体的完整接口。

\lecturefigure{slide-42-thanks-results.jpg}{结语：Cicero 的成绩、论文、代码与开放研究议程}{官方录像恢复幻灯片 V042；Stanford CS25 Lecture 14}

\teachervoice{在最后问答中，讲者说自己更关心 general planning，而不是为每个新任务单独 fine-tune 一个求解器。CoT 令人印象深刻，因为它跨任务通用，但它的规划能力仍有限；课堂判断是，更强、更一般的推理时方法应该存在。}

\begin{importantbox}{本讲留下的开放问题}
能否让同一个 foundation model 在不同任务中学习何时展开、展开什么、如何验证、何时停止，并把 search 产生的中间状态与外部世界可靠对齐？这比“多生成一些 token”更严格，也比为每个游戏手写求解器更通用。
\end{importantbox}

\subsection{本章小结}

对话案例证明 Cicero 的语言由可执行策略支撑；局限则说明它没有解决长期 intent、信任建模和通用 planning。系统的诚实倾向主要来自 evaluator 能力边界下的工程约束，而不是通用社会推理已经完成。

